\section{Related Works}

\label{sec:related-works}

\subsection{Learned sparse retrieval}
Learned sparse retrieval replaces fixed lexical statistics, such as BM25, with neural models that both reweight observed terms and expand the representation with semantically related terms absent from the original text. SPLADE-based methods demonstrate that contextual token weighting and vocabulary expansion can make sparse retrieval competitive with strong dense baselines~\citep{spladev1, spladev2, spladev3}. Recent work extends this line of research using LLM backbones and stronger training recipes~\citep{laconic}. While related efforts either adapt decoder backbones toward bidirectional behavior for retrieval~\citep{laconic} or jointly train dense and sparse retrieval within one model~\citep{bgem3,sparsejoint}, our work differs by studying sparse retrieval natively under causal attention within a unified backbone. 

\subsection{Multimodal Retrieval}
Recently, Multimodal Large Language Models (MLLMs)~\citep{gemini2, qwen3vl2025, gpt5} have demonstrated exceptional performance across a wide range of tasks. Consequently, universal multimodal embedding models~\citep{ vlm2vec, thinkthenembed2025, qwenvlembedding} have been built upon these MLLM backbones, achieving highly competitive results on benchmarks such as \mmebv~\citep{vlm2vecv2}. However, multimodal sparse retrieval remains much less explored. Existing LSR work typically introduces auxiliary modules or an extra model to interpret non-text modalities~\citep{visualsparta, stair, blipsparse, frozen, splare}. Such auxiliary modules restrict the sparse model to a specific modality pair and complicate its extension to new modalities. Our work addresses this gap by combining decoder-native sparse retrieval, unified dense and sparse modeling, and multimodal support within a single causal backbone. 